% Especificação REST embutida — sem referência externa

\subsection{Gateway central (\texttt{api-gateway})}

O gateway constitui o ponto único de entrada do framework na porta 3000. Todas as requisições de processamento de programas MiniPar transitam por ele, o que preserva o princípio de inversão de controle e concentra a orquestração em um módulo invariante. O endpoint principal aceita o código-fonte, a variante de tradução desejada e o modo de execução, devolvendo a árvore sintática abstrata serializada, a tabela de símbolos, o texto de saída, o código gerado quando aplicável, o registro dos passos do pipeline e, no modo distribuído, os resultados agregados dos \emph{workers}.

\begin{lstlisting}[caption={Corpo da requisição \texttt{POST /api/v1/process}}]
{
  "sourceCode": "class Main { void run() { println(\"ok\"); } }",
  "targetVariability": "INTERPRETER",
  "executionMode": "LOCAL"
}
\end{lstlisting}

Os valores admitidos para \texttt{targetVariability} compreendem \texttt{INTERPRETER}, \texttt{C}, \texttt{CPP}, \texttt{RUST}, \texttt{ASSEMBLY} e \texttt{PYTHON}. O campo \texttt{executionMode} aceita \texttt{LOCAL} ou \texttt{DISTRIBUTED\_SOCKETS}. A resposta bem-sucedida inclui os campos \texttt{success}, \texttt{historyId}, \texttt{output}, \texttt{ast}, \texttt{symbolTable}, \texttt{generatedCode}, \texttt{pipelineSteps} e, quando pertinente, \texttt{distributedResults}.

\begin{lstlisting}[caption={Trecho representativo da resposta agregada pelo gateway}]
{
  "success": true,
  "targetVariability": "C",
  "executionMode": "LOCAL",
  "output": "Compiled with gcc -O2\nok\n",
  "pipelineSteps": [
    "ms-front-end: parse",
    "ms-semantic: analyze",
    "ms-codegen-c: generate"
  ],
  "generatedCode": "/* program.c emitido pelo SimpleCCodeGenerator */"
}
\end{lstlisting}

O endpoint \texttt{GET /api/v1/variants} expõe as variantes registradas no vetor central de back-ends, permitindo que a interface de seleção de variabilidade se atualize sem alteração de código quando uma nova instância é adicionada. O endpoint \texttt{GET /api/v1/recommendations} consulta o histórico de compilações para sugerir configurações frequentes. O endpoint \texttt{GET /api/v1/services/health} agrega o estado de saúde de cada microsserviço configurado, consultando individualmente o recurso \texttt{/health} de cada contêiner.

\subsection{Microsserviço de análise léxica e sintática (\texttt{ms-front-end})}

Este ativo reutilizável executa a primeira fase do pipeline. Recebe o texto do programa na linguagem MiniPar 2026.1 e devolve a árvore sintática abstrata em formato JSON ou uma lista estruturada de erros sintáticos. O contrato permanece independente do back-end de tradução, o que materializa o desacoplamento exigido pela linha de produto de software.

\begin{lstlisting}[caption={Requisição e resposta de \texttt{POST /parse} (porta 3001)}]
// Requisicao
{ "sourceCode": "class Animal { string name; }" }

// Resposta sem erro
{
  "ast": { "type": "Program", "declarations": [ /* ... */ ] },
  "errors": []
}

// Resposta com erro sintático
{
  "ast": null,
  "errors": [{ "message": "Parser error at 2:11: Expected LBRACE" }]
}
\end{lstlisting}

\subsection{Microsserviço de análise semântica (\texttt{ms-semantic})}

O analisador semântico consome exclusivamente a árvore serializada produzida na fase anterior. A função \texttt{analyze\_ast\_full} desserializa o documento JSON, percorre a árvore com o visitador \texttt{SemanticAnalyzer} e constrói a tabela de símbolos hierárquica com escopos aninhados. Erros de herança inválida, declaração duplicada de classes ou membros repetidos retornam no vetor \texttt{errors}, interrompendo o pipeline antes da invocação do back-end.

\begin{lstlisting}[caption={Requisição e resposta de \texttt{POST /analyze} (porta 3002)}]
// Requisicao
{ "ast": { "type": "Program", "declarations": [ /* ... */ ] } }

// Resposta
{
  "ast": { "type": "Program", "declarations": [ /* ... */ ] },
  "symbolTable": {
    "scopes": [
      { "name": "global", "symbols": [ /* classes, canais, funções */ ] }
    ]
  },
  "errors": []
}
\end{lstlisting}

\subsection{Microsserviço interpretador (\texttt{ms-interpreter})}

A variante de execução direta sobre a árvore sintática abstrata implementa o ponto de adaptação \texttt{Interpreter.emit} dentro de um contêiner autônomo. O endpoint \texttt{POST /execute} recebe a árvore validada e opcionalmente a tabela de símbolos, executa o interpretador orientado a objetos com suporte a blocos \texttt{par}, \texttt{seq} e canais por soquete TCP, e retorna o texto capturado na saída padrão simulada.

\begin{lstlisting}[caption={Contrato de \texttt{POST /execute} (porta 3003)}]
// Requisicao
{
  "ast": { "type": "Program", "declarations": [ /* ... */ ] },
  "symbolTable": { "scopes": [ /* ... */ ] },
  "executionMode": "LOCAL",
  "target": "INTERPRETER"
}

// Resposta
{ "output": "woof\n", "exitCode": 0 }
\end{lstlisting}

\subsection{Microsserviços de geração de código}

Os back-ends de compilação compartilham o mesmo contrato de entrada: árvore sintática abstrata validada e tabela de símbolos. A diferença reside no ponto de adaptação implementado em cada instância. O serviço \texttt{ms-codegen-c} (porta 3004) produz código em linguagem C, compila com \texttt{gcc -O2} e retorna saída e artefato gerado. O serviço \texttt{ms-codegen-rust} (porta 3005) emite código Rust e invoca \texttt{rustc} quando disponível. O serviço \texttt{ms-codegen-arm} (porta 3006) gera assembly ARMv7 a partir da representação intermediária. O serviço \texttt{ms-codegen-python} (porta 3008), configurado como extensão da linha de produto, demonstra a instanciação de nova variante sem alteração dos trechos invariantes do núcleo.

\begin{lstlisting}[caption={Contrato comum \texttt{POST /generate}}]
// Requisicao
{
  "ast": { "type": "Program", "declarations": [ /* ... */ ] },
  "symbolTable": { "scopes": [ /* ... */ ] }
}

// Resposta tipica (back-end C)
{
  "output": "Compiled with gcc -O2\nok\n",
  "generatedCode": "int main(void) { /* ... */ }",
  "exitCode": 0
}
\end{lstlisting}

O back-end C++ reutiliza o mesmo contêiner e serviço do back-end C, alterando apenas o compilador invocado para \texttt{g++} com o sinalizador \texttt{-std=c++17}, conforme a classe \texttt{CppBackend} derivada por herança.

\subsection{Coordenador de paralelismo (\texttt{ms-parallel-coord})}

Quando o modo de execução \texttt{DISTRIBUTED\_SOCKETS} é selecionado e o programa-fonte não utiliza canais de comunicação distribuídos, o gateway encaminha a coordenação ao serviço na porta 3007. Este componente dispara, em paralelo, requisições de execução aos três \emph{workers} socket nas portas 9001 a 9003, cada um executando um programa MiniPar independente (\texttt{worker\_quicksort.minipar}, \texttt{worker\_matrix.minipar}, \texttt{worker\_factorial.minipar}). O protocolo de transporte consiste no envio do comando textual \texttt{RUN} seguido da leitura de um documento JSON com os campos \texttt{role}, \texttt{data}, \texttt{machine}, \texttt{ip} e \texttt{port}.

\begin{lstlisting}[caption={Requisição e resposta de \texttt{POST /coordinate}}]
// Requisicao (hosts opcionais; padrão: três workers do Compose)
{
  "executionMode": "DISTRIBUTED_SOCKETS",
  "ast": { "type": "Program", "declarations": [] }
}

// Resposta
{
  "output": "=== Coordenador Distribuido ===\n[PC1] QuickSort: ...\n...",
  "results": [
    { "role": "quicksort", "machine": "PC1", "data": "...", "ip": "172.x.x.x", "port": 9001 },
    { "role": "matrix", "machine": "PC2", "data": "...", "ip": "172.x.x.x", "port": 9002 },
    { "role": "factorial", "machine": "PC3", "data": "...", "ip": "172.x.x.x", "port": 9003 }
  ]
}
\end{lstlisting}

Quando o programa-fonte declara \texttt{c\_channel} e operações \texttt{send} ou \texttt{receive}, o interpretador estabelece comunicação direta entre processos por meio do módulo \texttt{socket\_channel.py}, dispensando o coordenador legado. O programa \texttt{14\_distributed\_menu.minipar} exemplifica o menu central que dispara os três \emph{workers}; o programa \texttt{15\_channels.minipar} valida a passagem de valor inteiro entre processos locais com saída esperada \texttt{42}.

\subsection{Tratamento de erros e encadeamento do pipeline}

O gateway interrompe o encadeamento na primeira fase que retorna erro. Erros sintáticos originam-se em \texttt{ms-front-end} com \texttt{ast} nulo; erros semânticos em \texttt{ms-semantic} com vetor \texttt{errors} não vazio; falhas de compilação externa propagam-se a partir do campo \texttt{exitCode} distinto de zero no \texttt{TranslationResult}. Cada passo bem-sucedido registra-se em \texttt{pipelineSteps}, o que permite auditoria da execução e reprodução dos testes documentados na Seção de validação.
